\documentclass[11pt,a4paper]{article}

%  XeLaTeX
\usepackage{fontspec}

\usepackage{algorithm}
\usepackage{algorithmic}
\usepackage{amsmath,amssymb,amsthm}
\usepackage{bm}
\usepackage[useregional]{datetime2}
\usepackage{enumitem}
\usepackage{booktabs}
\usepackage{float}
\usepackage{graphicx}
\usepackage{tabularray}
\usepackage{framed}
\usepackage[a4paper, margin=2.5cm]{geometry}
\usepackage{eso-pic}
\usepackage{mathtools}
\usepackage{tikz}
\usetikzlibrary{calc}
\usepackage{xcolor}

\usepackage[backend=biber, style=apa, sorting=nyt]{biblatex}
\addbibresource{D:/github/ENEL445/report/references.bib}
\graphicspath{{D:/github/ENEL445/figs/}}

\usepackage[colorlinks=true, linkcolor=blue, citecolor=blue, urlcolor=blue]{hyperref}
\hypersetup{
  pdfauthor   = {David Ewing},
  pdfsubject  = {\jobname.tex},
  pdfkeywords = {source: \jobname.tex}
}

% Running margin label
\newcommand{\mymarginlabel}{\textcopyright\ 2026 David Ewing dew59@uclive.ac.nz | \DTMnow\ | \jobname.tex}
\AddToShipoutPictureBG{%
  \begin{tikzpicture}[remember picture,overlay]
    \node[rotate=90, anchor=north]
      at ($(current page.north west)+(1.4cm,-13cm)$)
      {\scriptsize\ttfamily \mymarginlabel};
  \end{tikzpicture}%
}

% Custom commands
\newcommand{\E}{\mathbb{E}}
\newcommand{\R}{\mathbb{R}}
\newcommand{\N}{\mathcal{N}}
\newcommand{\ELBO}{\text{ELBO}}
\newcommand{\tr}{\text{tr}}
\setlength{\parindent}{0pt}

\title{
Case Study 4 --- Hierarchical Bayesian Linear Regression:\\[0.4em]
Gradient-Based Optimisation and Posterior Variance Collapse
}
\author{
  David Ewing (82171165)
}
\date{11 May 2026}

\begin{document}

\maketitle

\begin{abstract}
\setlength{\parindent}{0pt}
\setlength{\parskip}{0.7em}

This report applies the gradient-based optimisation methods from ENEL445 to variational
inference in hierarchical Bayesian linear regression.
Extending the simple Gaussian regression setting, this model introduces $J=5$ group-level
random effects $u_j \sim \mathcal{N}(0,\tau_u^{-1})$ alongside the noise precision
$\tau_e$ as additional latent variables with Gamma priors.
The Evidence Lower Bound (ELBO) for this conjugate model admits closed-form CAVI updates,
but also serves as a smooth objective for gradient-based methods, with $D=19$ unconstrained
parameters.

Four optimisation methods are applied to maximise the ELBO:
Coordinate Ascent Variational Inference (CAVI), gradient ascent with Armijo backtracking,
Newton's method with regularised finite-difference Hessian, and BFGS with Wolfe conditions.
The variational family is a mean-field product
$q(\bm{\beta})\prod_j q(u_j) \cdot q(\tau_e)\cdot q(\tau_u)$
with Gaussian random effects and Gamma precision approximations, and a
Cholesky reparameterisation for $\bm{\Sigma}_\beta$.

The reference posterior is obtained from a conjugate blocked Gibbs sampler
using 3 chains of 5000 iterations (500 burn-in each).
A striking finding is severe posterior variance collapse for $\beta_0$
(SD ratio $\approx 0.13$ for CAVI and BFGS), driven by the mean-field factorisation
absorbing shared group-level variance into the random effects $u_j$ rather than
the population intercept.
Newton's method fails to converge to the correct optimum, highlighting the
sensitivity of second-order methods to poor Hessian conditioning in high-dimensional
hierarchical models.

\end{abstract}

\newpage

%─────────────────────────────────────────────────────────────────────────────
\section{Introduction}
%─────────────────────────────────────────────────────────────────────────────

Hierarchical Bayesian linear regression is the natural extension of the simple
Gaussian regression baseline (Case Study~1) to settings with group structure.
It provides the first non-trivial test of whether gradient-based VI methods can
navigate a substantially higher-dimensional posterior whilst competing with the conjugate
Gibbs reference, and whether \emph{posterior variance collapse} --- a known failure
mode of mean-field VI --- manifests and can be detected.

Bayesian inference provides principled uncertainty quantification but requires computing
intractable posterior distributions. For most real-world models, the posterior integral,
\begin{equation*}
    p(\bm{\theta} \mid \bm{y}) = \frac{p(\bm{y} \mid \bm{\theta})\,p(\bm{\theta})}
    {\displaystyle\int p(\bm{y} \mid \bm{\theta})\,p(\bm{\theta})\, d\bm{\theta}},
\end{equation*}
is intractable.
In this case study,
$\bm{\theta} = (\bm{\beta},\,\mathbf{u},\,\tau_e,\,\tau_u)$
--- the population-level regression coefficients, the $J$ group-level random effects,
the observation noise precision, and the random-effect precision --- so the posterior
of interest is
\begin{equation*}
    p(\bm{\beta},\mathbf{u},\tau_e,\tau_u \mid \mathbf{y})
    \propto p(\mathbf{y}\mid\bm{\beta},\mathbf{u},\tau_e)\,
             p(\bm{\beta})\,p(\mathbf{u}\mid\tau_u)\,
             p(\tau_e)\,p(\tau_u).
\end{equation*}
The substitution $\bm{\theta} = (\bm{\beta},\mathbf{u},\tau_e,\tau_u)$ is specific to
this case study; other case studies define $\bm{\theta}$ according to their own model
parameters.
Unlike Case Study~1, this posterior does not factorise into a product of independent
conjugate terms; the random effects $u_j$ and the precision $\tau_u$ are tightly coupled,
making direct computation challenging even for Gaussian likelihoods.

Variational Inference (VI) converts this integration into optimisation by finding
a tractable approximate posterior $q$ that maximises the Evidence Lower Bound (ELBO).
This report investigates how the gradient-based methods from ENEL445 perform on this
higher-dimensional problem, compared to the conjugate blocked Gibbs sampler as the
reference.

A well-known failure mode of mean-field VI is \emph{posterior variance collapse}:
the optimised approximation $q(\bm{\theta})$ is systematically overconfident,
underestimating posterior uncertainty.
This arises because the mean-field factorisation discards posterior \emph{correlations}
between blocks --- here, between the intercept $\beta_0$ and the group means $u_j$
--- forcing each factor to concentrate more tightly than the true marginal.%
\footnote{See Section~\ref{sec:variance-collapse} for the hierarchical mechanism;
the general derivation is in Case Study~1 (Appendix~A).}

This report applies four optimisation methods --- CAVI, gradient ascent, Newton's
method, and BFGS --- to the ELBO and compares their convergence behaviour,
computational cost, and posterior quality.
Of particular interest is whether the SD ratio
$\sigma_{\mathrm{VI}}/\sigma_{\mathrm{Gibbs}}$ departs substantially from unity,
signalling collapse, and which parameters are most affected.

Throughout all gradient-based methods, the variational covariance $\bm{\Sigma}_\beta$
must remain symmetric positive definite at every iteration.%
\footnote{This structural constraint is satisfied automatically via Cholesky
reparameterisation; derivation in Case Study~1 (Appendix~B).}

The classical foundation for the variational framework used here is least-squares
regression, as developed in lectures by Prof.\ Roy S.\ Smith (Erskine Visiting Fellow,
ETH~Z\"{u}rich, ENEL445 2026).
The design matrix $\bm{X}$%
\footnote{Denoted $\Phi$ and called the \emph{regressor matrix} in lectures
(Prof.\ R.\,S.\ Smith, ENEL445 2026).}
and coefficient vector $\bm{\beta}$%
\footnote{Denoted $\theta$ in lectures.}
follow the standard over-determined formulation.
The noise model adopted here --- errors on $\bm{y}$ only, with noise precision
$\tau_e$ --- corresponds to the equation-error form; this is an explicit modelling
assumption rather than a derived result.

\subsection{Bayesian Regression Model}

The hierarchical Bayesian linear regression model is
\begin{align}
y_{ij} \mid \bm{\beta}, u_j, \tau_e
  &\sim \mathcal{N}\!\bigl(\mathbf{x}_{ij}^T\bm{\beta} + u_j,\; \tau_e^{-1}\bigr),
  \quad i=1,\ldots,n_j,\; j=1,\ldots,J, \\
u_j \mid \tau_u &\sim \mathcal{N}(0,\; \tau_u^{-1}), \\
\bm{\beta} &\sim \N(\bm{0},\; \lambda^{-1}\bm{I}_p), \quad \lambda=0.1, \\
\tau_e &\sim \mathrm{Gamma}(\alpha_e,\; \gamma_e), \quad \alpha_e=\gamma_e=1, \\
\tau_u &\sim \mathrm{Gamma}(\alpha_u,\; \gamma_u), \quad \alpha_u=\gamma_u=1.
\end{align}
Here $\bm{X}$ is the design matrix with rows $\mathbf{x}_{ij}^T = (1, x_{ij})$,
$\mathbf{y}$ is the response vector, $\bm{\beta}$ is the fixed-effect coefficient
vector, $u_j$ is the group-$j$ random effect, and $\tau_e$, $\tau_u$ are the
noise and random-effect precisions respectively.

\subsection{Mean-Field Variational Approximation}

The mean-field variational family factorises as
\begin{align}
q(\bm{\beta},\mathbf{u},\tau_e,\tau_u)
  &= q(\bm{\beta})\;\prod_{j=1}^{J} q(u_j)\;\cdot\; q(\tau_e)\;\cdot\; q(\tau_u),
\end{align}
with
\begin{align}
q(\bm{\beta}) &= \N(\bm{\mu}_\beta,\,\bm{\Sigma}_\beta), \\
q(u_j) &= \mathcal{N}(m_j,\, v_j), \quad j=1,\ldots,J, \\
q(\tau_e) &= \mathrm{Gamma}(a_e,\, b_e), \\
q(\tau_u) &= \mathrm{Gamma}(a_u,\, b_u).
\end{align}
The variational parameters are
$\bm{\theta} = (\bm{\mu}_\beta, \bm{L}, m_1,\ldots,m_J, \log v_1,\ldots,\log v_J,
\log a_e, \log b_e, \log a_u, \log b_u)$,
where $\bm{\Sigma}_\beta = \bm{L}\bm{L}^T$ and $\bm{L}$ is lower-triangular with
positive diagonal, giving $D = p + p(p+1)/2 + J + J + 4 = 2 + 3 + 5 + 5 + 4 = 19$
unconstrained parameters for $p=2$, $J=5$.

\subsection{Key Notation}

\begin{center}
\begin{tabular}{@{}ll@{}}
\toprule
Symbol & Meaning \\
\midrule
\multicolumn{2}{@{}l}{\textbf{Bayesian model}} \\
$\bm{X}$ & Design matrix ($N \times p$), rows $\mathbf{x}_{ij}^T$ \\
$\mathbf{y}$ & Continuous response vector \\
$\bm{\beta}$ & Fixed-effect coefficient vector \\
$u_j$ & Group random effect for group $j$ \\
$\tau_e, \tau_u$ & Noise and random-effect precision \\
$N, J, p$ & Total observations, groups, predictors \\
\midrule
\multicolumn{2}{@{}l}{\textbf{Variational approximation}} \\
$q(\bm{\beta})$ & Variational Gaussian for fixed effects \\
$\bm{\mu}_\beta, \bm{\Sigma}_\beta$ & Mean and covariance of $q(\bm{\beta})$ \\
$\bm{L}$ & Cholesky factor, $\bm{\Sigma}_\beta=\bm{L}\bm{L}^T$ \\
$m_j, v_j$ & Mean and variance of $q(u_j)$ \\
$a_e, b_e$ & Gamma parameters of $q(\tau_e)$ \\
$a_u, b_u$ & Gamma parameters of $q(\tau_u)$ \\
$\mathcal{L}(\bm{\theta})$ & ELBO (maximised) \\
\midrule
\multicolumn{2}{@{}l}{\textbf{Reparameterisation}} \\
$\bm{L}$ & Lower-triangular Cholesky factor, $\bm{\Sigma}_\beta=\bm{L}\bm{L}^T$ \\
$\tilde{\ell}_{ii}$ & Log-diagonal entry, $L_{ii}=e^{\tilde{\ell}_{ii}}$ (ensures $L_{ii}>0$) \\
$\log v_j$ & Log-variance, $v_j=e^{\log v_j}>0$ \\
$\bm{\theta}$ & Unconstrained optimisation vector ($D=19$) \\
\midrule
\multicolumn{2}{@{}l}{\textbf{Optimisation}} \\
$\bm{d}^{(t)}$ & Search direction at iteration $t$ \\
$\alpha_t$ & Step size (line-search parameter) at iteration $t$ \\
$\nabla\mathcal{L}$ & Gradient of the ELBO w.r.t.\ $\bm{\theta}$ \\
$\nabla^2\mathcal{L}$ & Hessian of the ELBO w.r.t.\ $\bm{\theta}$ \\
$\bm{B}_t$ & BFGS inverse-Hessian approximation at iteration $t$ \\
\midrule
\multicolumn{2}{@{}l}{\textbf{Special functions}} \\
$\psi(x)$ & Digamma function, $\psi(x)=\tfrac{d}{dx}\log\Gamma(x)$ \\
$\psi^{(1)}(x)$ & Trigamma function, $\psi^{(1)}(x)=\tfrac{d}{dx}\psi(x)$ \\
\midrule
\multicolumn{2}{@{}l}{\textbf{Lecture correspondence} (Prof.\ R.\,S.\ Smith, ENEL445 2026)} \\
$J(\theta)$ & Cost function $=\mathcal{L}(\bm{\theta})$ in this report \\
$\Phi$ & Regressor matrix $=\bm{X}$ in this report \\
\bottomrule
\end{tabular}
\end{center}

\subsection{Numerical Study and Required Optimisation Components}

This report formulates variational inference for hierarchical Bayesian linear regression
as an optimisation problem and applies four methods: Coordinate Ascent Variational
Inference (CAVI), gradient ascent, Newton's method, and BFGS. Each method is
characterised by four components from the course framework: entry conditions, search
direction, line search, and exit conditions. All four methods share the same initial
point so that comparisons reflect method behaviour rather than initialisation.

\subsubsection*{Entry Conditions}

All methods are initialised at the prior mean:
$\bm{\mu}_\beta = \bm{0}$, $\bm{\Sigma}_\beta = \bm{I}_p$,
$m_j = 0$, $v_j = 1$ for $j=1,\ldots,J$,
$a_e = \alpha_e$, $b_e = \gamma_e$, $a_u = \alpha_u$, $b_u = \gamma_u$.
In the unconstrained parameterisation this corresponds to $\bm{L} = \bm{I}_p$
(so $\tilde{\ell}_{ii} = 0$), $\log v_j = 0$, and the log-space Gamma parameters
at their prior values.
CAVI uses the same starting variational parameters.

\subsubsection*{Search Direction}

For all gradient-based methods, the proposed search direction must be an ascent
direction before a step is accepted:
\begin{align}
\nabla \ELBO(\bm{\theta}^{(t)})^T \bm{d}^{(t)} > 0.
\end{align}
Gradient ascent uses $\bm{d}^{(t)} = \nabla\mathcal{L}(\bm{\theta}^{(t)})$ directly.
Newton's method uses $\bm{d}^{(t)} = -[\nabla^2\mathcal{L}(\bm{\theta}^{(t)})]^{-1}\nabla\mathcal{L}(\bm{\theta}^{(t)})$,
falling back to the gradient direction if the Hessian is not positive definite.
BFGS approximates the inverse Hessian as $\bm{B}_t$ and uses
$\bm{d}^{(t)} = \bm{B}_t \nabla\mathcal{L}(\bm{\theta}^{(t)})$.
CAVI replaces the search direction with closed-form coordinate updates derived
analytically from the model.

\subsubsection*{Line Search}

The step size $\alpha^{(t)}$ is determined per method.
Gradient ascent uses Armijo backtracking: starting from an initial step size,
$\alpha^{(t)}$ is halved until sufficient increase in the ELBO is achieved.
BFGS uses Wolfe conditions, combining sufficient increase with a curvature condition
that ensures
\begin{align}
\bm{y}_t^T \bm{s}_t > 0,
\end{align}
where $\bm{s}_t = \bm{\theta}^{(t+1)} - \bm{\theta}^{(t)}$ and
$\bm{y}_t = \nabla\mathcal{L}(\bm{\theta}^{(t+1)}) - \nabla\mathcal{L}(\bm{\theta}^{(t)})$.
This positivity condition is required for the BFGS update to preserve positive
definiteness of $\bm{B}_t$.
Newton's method uses a fixed step size of 1 unless the Hessian is indefinite.
CAVI has no line search: each coordinate update exactly maximises the ELBO with
respect to that parameter in closed form.

\subsubsection*{Exit Conditions}

Exit conditions are applied uniformly across the optimisation methods.
Iteration terminates when any of the following conditions is satisfied:
\begin{enumerate}[label=(\roman*)]
  \item Gradient norm: $\lVert \nabla \ELBO \rVert < \varepsilon_g$.
  \item Relative ELBO change: $\lvert \Delta \ELBO \rvert / \lvert \ELBO \rvert < \varepsilon_r$.
  \item Parameter change: $\lVert \Delta \bm{\theta} \rVert < \varepsilon_\theta$.
  \item Maximum iterations: $t \geq t_{\max}$.
\end{enumerate}
CAVI additionally monitors the maximum relative change in its variational parameters:
\begin{align}
\max_i \frac{\lvert \theta_i^{(t+1)} - \theta_i^{(t)} \rvert}{\lvert \theta_i^{(t+1)} \rvert} < \varepsilon.
\end{align}

%─────────────────────────────────────────────────────────────────────────────
\section{Model Formulation}
%─────────────────────────────────────────────────────────────────────────────

\subsection{Hierarchical Bayesian Linear Regression}

The model is
\begin{align}
  y_{ij} \mid \bm{\beta}, u_j, \tau_e
    &\sim \mathcal{N}\!\bigl(\mathbf{x}_{ij}^T\bm{\beta} + u_j,\; \tau_e^{-1}\bigr),
    \quad i=1,\ldots,n_j,\; j=1,\ldots,J,\\
  u_j \mid \tau_u &\sim \mathcal{N}(0,\; \tau_u^{-1}),\\
  \bm{\beta} &\sim \N(\bm{0},\; \lambda^{-1}\bm{I}_p), \quad \lambda=0.1,\\
  \tau_e &\sim \mathrm{Gamma}(\alpha_e,\; \gamma_e), \quad \alpha_e=\gamma_e=1,\\
  \tau_u &\sim \mathrm{Gamma}(\alpha_u,\; \gamma_u), \quad \alpha_u=\gamma_u=1,
\end{align}
where $\mathbf{x}_{ij} = (1, x_{ij})^T \in \R^2$ and $\sigma(\cdot)$ is the logistic sigmoid
(not used here; included for consistency with notation across case studies).

The test case uses $N=150$ observations in $J=5$ groups (30 per group),
$x_{ij} \sim \mathrm{Uniform}(-2,2)$, and true parameters
$\bm{\beta}_{\mathrm{true}} = (1.0,\; 2.0)$,
$\tau_{e,\mathrm{true}}=4.0$ ($\sigma_e = 0.5$),
$\tau_{u,\mathrm{true}}=2.0$ ($\sigma_u \approx 0.707$).

\subsection{Mean-Field Variational Approximation}

The mean-field variational family factorises as
\begin{align}
  q(\bm{\beta},\mathbf{u},\tau_e,\tau_u)
  &= q(\bm{\beta})\;\prod_{j=1}^{J} q(u_j)\;\cdot\; q(\tau_e)\;\cdot\; q(\tau_u),
\end{align}
with
\begin{align}
  q(\bm{\beta}) &= \N(\bm{\mu}_\beta,\,\bm{\Sigma}_\beta),\\
  q(u_j) &= \mathcal{N}(m_j,\, v_j), \quad j=1,\ldots,J,\\
  q(\tau_e) &= \mathrm{Gamma}(a_e,\, b_e),\\
  q(\tau_u) &= \mathrm{Gamma}(a_u,\, b_u).
\end{align}
The covariance $\bm{\Sigma}_\beta = \bm{L}\bm{L}^T$ is parameterised by its Cholesky
factor $\bm{L}$ to ensure positive definiteness throughout optimisation.

\subsection{Unconstrained Parameterisation and ELBO}

The unconstrained parameter vector $\bm{\theta} \in \R^{19}$ is
\begin{equation}
  \bm{\theta} = \bigl(\bm{\mu}_\beta,\;
    \tilde{\ell}_{11}, \ell_{21}, \tilde{\ell}_{22},\;
    m_1,\ldots,m_5,\;
    \log v_1,\ldots,\log v_5,\;
    \log a_e, \log b_e, \log a_u, \log b_u
  \bigr),
\end{equation}
where $L_{ii} = e^{\tilde{\ell}_{ii}} > 0$ and $v_j = e^{\log v_j} > 0$.

The ELBO is
\begin{align}
  \mathcal{L}(\bm{\theta})
  &= \E_q\!\left[\log p(\mathbf{y}\mid\bm{\beta},\mathbf{u},\tau_e)\right]
    + \E_q\!\left[\log p(\mathbf{u}\mid\tau_u)\right]
    \nonumber\\
  &\quad - \mathrm{KL}\bigl(q(\bm{\beta})\|p(\bm{\beta})\bigr)
    + \E_q\!\left[\log p(\tau_e)\right] + H[q(\tau_e)]
    \nonumber\\
  &\quad + \E_q\!\left[\log p(\tau_u)\right] + H[q(\tau_u)]
    + \sum_j H[q(u_j)],
\end{align}
where all terms are available in closed form for Gaussian--Gamma conjugacy.

\subsection{CAVI Update Equations}

The coordinate ascent updates cycle over:
\begin{align}
  \bm{\Sigma}_\beta^{-1} &\leftarrow \lambda\bm{I} + \E[\tau_e]\,\bm{X}^T\bm{X}, \label{eq:cavi_beta_cov}\\
  \bm{\mu}_\beta &\leftarrow \bm{\Sigma}_\beta\,\E[\tau_e]\,\bm{X}^T(\mathbf{y} - \bm{m}[\mathrm{groups}]),\\
  v_j^{-1} &\leftarrow \E[\tau_e]\,n_j + \E[\tau_u], \quad
  m_j \leftarrow v_j\,\E[\tau_e]\!\sum_{i\in j}\!\bigl(y_{ij} - \mathbf{x}_{ij}^T\bm{\mu}_\beta\bigr),\\
  a_e &\leftarrow \alpha_e + \tfrac{N}{2}, \quad
  b_e \leftarrow \gamma_e + \tfrac{1}{2}\E\!\left[\|\mathbf{y} - \bm{X}\bm{\beta} - \bm{m}[\mathrm{groups}]\|^2\right],\\
  a_u &\leftarrow \alpha_u + \tfrac{J}{2}, \quad
  b_u \leftarrow \gamma_u + \tfrac{1}{2}\sum_j\!(m_j^2 + v_j),
\end{align}
where $\E[\tau_e] = a_e/b_e$ and $\E[\tau_u] = a_u/b_u$.

%─────────────────────────────────────────────────────────────────────────────
\section{Conjugate Gibbs Sampler}
%─────────────────────────────────────────────────────────────────────────────

The hierarchical Gaussian model is fully conjugate, so a blocked Gibbs sampler is
available with closed-form full conditionals.
Each sweep samples:
\begin{align}
  \bm{\beta} \mid \mathbf{u}, \tau_e, \mathbf{y}
    &\sim \mathcal{N}(V_\beta\,\bm{X}^T\tau_e(\mathbf{y}-\bm{m}[\mathrm{groups}]),\; V_\beta),
    \quad V_\beta = (\lambda\bm{I} + \tau_e\bm{X}^T\bm{X})^{-1},\\
  u_j \mid \bm{\beta}, \tau_e, \tau_u, \mathbf{y}
    &\sim \mathcal{N}(V_j\tau_e\!\sum_{i\in j}(y_{ij} - \mathbf{x}_{ij}^T\bm{\beta}),\; V_j),
    \quad V_j = (\tau_e n_j + \tau_u)^{-1},\\
  \tau_e \mid \bm{\beta}, \mathbf{u}, \mathbf{y}
    &\sim \mathrm{Gamma}\!\left(\alpha_e + \tfrac{N}{2},\;
      \gamma_e + \tfrac{1}{2}\|\mathbf{y}-\bm{X}\bm{\beta}-\mathbf{u}[\mathrm{groups}]\|^2\right),\\
  \tau_u \mid \mathbf{u}
    &\sim \mathrm{Gamma}\!\left(\alpha_u + \tfrac{J}{2},\;
      \gamma_u + \tfrac{1}{2}\sum_j u_j^2\right).
\end{align}
The sampler runs for $N_{\mathrm{warmup}} = 500$ burn-in and $N = 5000$ retained
iterations per chain, with 3 independent chains.

%─────────────────────────────────────────────────────────────────────────────
\section{Posterior Variance Collapse}
\label{sec:variance-collapse}
%─────────────────────────────────────────────────────────────────────────────

Posterior variance collapse means that the variational covariance $\bm{\Sigma}_\beta$
understates the true uncertainty in the regression coefficients.
The collapsed posterior then looks more precise than the data justify.
In this hierarchical model the effect is more severe than in simple linear regression
(Case Study~1) because the hierarchical coupling creates strong posterior correlations
that the mean-field factorisation is forced to discard.

\subsection{Why Collapse Occurs}

The mean-field approximation factorises the variational family as
\begin{align}
  q(\bm{\beta},\mathbf{u},\tau_e,\tau_u)
  = q(\bm{\beta})\;\prod_{j=1}^{J} q(u_j)\;\cdot\; q(\tau_e)\;\cdot\; q(\tau_u).
\end{align}
This factorisation is convenient, but it discards all posterior covariances
between the global parameters $\bm{\beta}$ and the group-level random effects
$u_1,\ldots,u_J$.

In the true posterior these two blocks are strongly negatively correlated.
Given the model $y_{ij} = \beta_0 + \beta_1 x_{ij} + u_j + \varepsilon_{ij}$,
the data constrain the \emph{group mean} $\beta_0 + u_j$ far more tightly
than either $\beta_0$ or $u_j$ alone.
If a posterior sample has $\beta_0$ large, the group effects $u_j$ must be
correspondingly small, and vice versa.
The mean-field assumption severs this link: $q(\beta_0)$ and each $q(u_j)$ are
forced to be independent, so the variational factor $q(\beta_0)$ is fitted to
the marginal likelihood of $\beta_0$ \emph{as if the group effects were fixed
at their posterior means}.
This narrows $q(\beta_0)$ relative to the true marginal posterior, producing
the collapse.

The slope $\beta_1$ is not confounded with the random effects in the same way:
the predictor $x_{ij}$ varies within each group and is not absorbed by $u_j$,
so $q(\beta_1)$ is much less affected.

\paragraph{Direction of the KL divergence.}
Standard VI minimises $\mathrm{KL}[q\|p]$, which penalises $q$ heavily for
putting mass where $p$ is small, but is lenient when $q$ fails to cover parts
of the posterior tails.
The practical effect is that the variational approximation sits on the
high-density centre of $p(\bm{\theta}\mid\mathbf{y})$ rather than covering
the full posterior spread, reinforcing the narrowing in $\bm{\Sigma}_\beta$.

\subsection{Diagnostic}

For each scalar parameter, the posterior standard-deviation ratio
\begin{align}
  r = \frac{\sigma_{\mathrm{VI}}}{\sigma_{\mathrm{Gibbs}}}
\end{align}
is computed using the Gibbs sampler as the reference (since the exact posterior
is not available in closed form for a hierarchical model).
A ratio below one confirms collapse; a ratio near one indicates the VI
approximation matches the Gibbs marginal.

In this study, the well-converged VI methods (CAVI, gradient ascent, BFGS)
all produce $r \approx 0.13$ for $\beta_0$ --- the VI standard deviation is
roughly one-eighth of the Gibbs standard deviation, a severe collapse.
In contrast, $r \approx 1.0$ for $\beta_1$.
The intercept--random-effect confounding is thus the dominant source of
collapse, and the slope is effectively immune.

\subsection{Approaches Considered}

Three approaches to reducing or diagnosing collapse are considered.
Only entropy regularisation (Approach~3) is implemented; the first two are
discussed to motivate that choice.
The mechanistic derivation and gradient formulae are given in Case Study~1
(Appendix~A); only the hierarchical-specific adaptations are noted here.

\subsubsection{Informative Prior on $\tau_e$}

A stronger prior on the noise precision $\tau_e$ limits how large $\E[\tau_e]$
can become.
Since a large expected precision drives down $\bm{\Sigma}_\beta$ in the CAVI
update $\bm{\Sigma}_\beta^{-1} = \lambda\bm{I} + \E[\tau_e]\bm{X}^T\bm{X}$,
controlling $\tau_e$ can indirectly slow the collapse.
The limitation is that this changes the statistical model and requires a
defensible prior choice that may not generalise to other problems.

\subsubsection{Minimum Variance Constraint}

A direct lower bound $\Sigma_{\beta,jj} \geq \sigma^2_{\min}$ can be imposed
on each diagonal entry of $\bm{\Sigma}_\beta$.
This prevents the intercept variance from reaching pathologically small values
but introduces an external threshold that is not determined by the data.

\subsubsection{Entropy Regularisation}

The most direct objective-based approach is entropy regularisation:
\begin{align}
  \mathcal{L}_{\mathrm{reg}}(\bm{\theta})
  = \mathcal{L}(\bm{\theta}) + \lambda\, H[q(\bm{\beta})],
\end{align}
where $H[q(\bm{\beta})] = \tfrac{1}{2}\log\det(2\pi e\,\bm{\Sigma}_\beta)$
is the entropy of the variational Gaussian for $\bm{\beta}$.
The regularisation term rewards broader covariance estimates and directly
counters the collapse tendency.
For the Cholesky-parameterised covariance used here, the gradient with respect
to each Cholesky factor gains a closed-form correction (derivation in Case
Study~1, Appendix~A).

%─────────────────────────────────────────────────────────────────────────────
\section{Optimisation Problem Formulation}
%─────────────────────────────────────────────────────────────────────────────

\subsection*{Decision Variables}

Maximisation is carried out over $\bm{\theta} \in \R^{19}$ as defined above.
The constraints $\bm{\Sigma}_\beta \succ 0$ (Cholesky), $v_j > 0$ ($\log v_j$
unconstrained), $a_e, b_e, a_u, b_u > 0$ ($\log$-parameterised) are all
satisfied automatically.

\subsection*{Objective Function}

Maximise the ELBO $\mathcal{L}(\bm{\theta}):\R^{19}\to\R$.
The function is smooth and bounded above by $\log p(\mathbf{y})$.
For the conjugate Gaussian model, CAVI converges to the global optimum of the
mean-field ELBO; gradient-based methods may find the same optimum or converge
to a saddle point if the Hessian is ill-conditioned.

\subsection*{Optimisation Components}

\begin{description}
\item[CAVI.]
  Closed-form coordinate updates (Equations~\ref{eq:cavi_beta_cov}--above).
  Unit step size; exit when $|\Delta\mathcal{L}| < 10^{-8}$.

\item[Gradient ascent.]
  Direction: $\bm{d}^{(t)} = \nabla\mathcal{L}(\bm{\theta}^{(t)})$ via central FD
  ($h=10^{-5}$).
  Armijo backtracking from $\alpha_0=0.5$, $c_1=0.1$, $\rho=0.5$.
  Exit: $|\Delta\mathcal{L}| < 10^{-7}$ or 2000 iterations.

\item[Newton's method.]
  Direction: $\bm{d}^{(t)} = -\bm{H}_{\mathrm{reg}}^{-1}\bm{g}^{(t)}$ using a
  $19\times 19$ FD Hessian.
  Regularisation ensures the direction is an ascent direction.
  Armijo backtracking; exit at 50 iterations.

\item[BFGS.]
  Quasi-Newton with inverse-Hessian approximation updated via the BFGS rank-2 formula.
  Wolfe conditions via scipy; exit at $\|\bm{g}\|<10^{-7}$ or 2000 iterations.
\end{description}

%─────────────────────────────────────────────────────────────────────────────
\section{Numerical Study}
%─────────────────────────────────────────────────────────────────────────────

\subsection{Test Case Design}

\begin{itemize}
  \item Model: $y_{ij} \sim \mathcal{N}(\beta_0 + \beta_1 x_{ij} + u_j,\; \tau_e^{-1})$
  \item True parameters: $\bm{\beta}_{\mathrm{true}} = (1.0,\; 2.0)$,
    $\tau_{e,\mathrm{true}}=4.0$, $\tau_{u,\mathrm{true}}=2.0$
  \item Groups: $J=5$, $n_j=30$, $N=150$
  \item Covariates: $x_{ij} \sim \mathrm{Uniform}(-2, 2)$
  \item Prior: $\bm{\beta} \sim \N(\bm{0}, 10\bm{I})$;
    $\tau_e, \tau_u \sim \mathrm{Gamma}(1, 0.1)$
  \item Seed: \texttt{82171165} (student ID)
\end{itemize}

\subsection{Gradient Verification}

Before running any optimiser, the gradient of the ELBO with respect to the
19-dimensional unconstrained vector $\bm{\theta}$ was verified against central
finite differences at a CAVI warm-start point.
All 19 components agreed to relative error below $10^{-3}$.

\begin{figure}[H]
  \centering
  \includegraphics[width=0.95\linewidth]{hierarchical_linear_elbo_gradient_check.png}
  \caption{Finite-difference gradient verification at the CAVI warm-start:
    $h=10^{-5}$ versus $h=10^{-7}$.
    Bars coincide for all 19 components, confirming correct gradient implementation.}
  \label{fig:gradcheck}
\end{figure}

\subsection{ELBO Landscape}

Figure~\ref{fig:landscape} shows a 2-D slice of the ELBO surface along
$\mu_{\beta_0}$ and $\mu_{\beta_1}$ with all other parameters held at the warm start.
The surface is smooth with a clear maximum near the true values.

\begin{figure}[H]
  \centering
  \includegraphics[width=0.60\linewidth]{hierarchical_linear_elbo_landscape.png}
  \caption{ELBO landscape slice in the $(\mu_{\beta_0}, \mu_{\beta_1})$ plane.
    Red star: CAVI warm start; white cross: true $\bm{\beta}$.}
  \label{fig:landscape}
\end{figure}

%─────────────────────────────────────────────────────────────────────────────
\section{Results}
%─────────────────────────────────────────────────────────────────────────────

\subsection{Generated Data}

Figure~\ref{fig:data} shows the 150 generated observations stratified by group.
Group random effects $u_j$ visibly shift the regression lines upward or downward
relative to the population curve.

\begin{figure}[H]
  \centering
  \includegraphics[width=0.80\linewidth]{hierarchical_linear_data_scatter.png}
  \caption{Generated hierarchical data ($N=150$, $J=5$).
    Dotted lines show group-specific regressions; the solid line is the population mean.}
  \label{fig:data}
\end{figure}

\begin{figure}[H]
  \centering
  \includegraphics[width=0.92\linewidth]{hierarchical_linear_likelihood_posterior.png}
  \caption{Likelihood and CAVI posterior for the population-level regression.
    The red dashed line is the true population mean $y=1.0+2.0x$; the tomato band is the
    95\% noise interval $\pm 2\sigma_e = \pm 1.0$, showing where individual observations
    fall around the population mean (ignoring group shifts).
    The blue band is the CAVI posterior 95\% credible interval for the population
    regression mean; $N=150$ observations across $J=5$ groups pull the posterior
    slope close to the true value, but the intercept is visibly biased low
    ($\hat{\beta}_0 \approx 0.46$ versus $\beta_{0,\mathrm{true}} = 1.0$).
    This downward bias is a direct consequence of variance collapse: the CAVI intercept
    absorbs less of the group-mean signal than the true posterior marginal would.}
  \label{fig:likpost}
\end{figure}

\subsection{ELBO Convergence}

Figure~\ref{fig:elbocv} shows the ELBO history for all four methods.
CAVI and BFGS converge to essentially the same final ELBO value ($-132.08$).
Gradient ascent converges but requires the full 2000-iteration budget,
reflecting the slow convergence of first-order methods on a 19-dimensional problem.
Newton's method converges in only 50 iterations but to a markedly worse ELBO
($-165.28$), indicating failure to escape a saddle region due to ill-conditioned
Hessian in the high-dimensional parameter space.

\begin{figure}[H]
  \centering
  \includegraphics[width=\linewidth]{hierarchical_linear_elbo_convergence.png}
  \caption{ELBO convergence for CAVI (left) and gradient-based methods (right).
    Newton's method converges to a suboptimal point; CAVI and BFGS agree.}
  \label{fig:elbocv}
\end{figure}

\subsection{Step-Size and Condition-Number Behaviour}

Figure~\ref{fig:stepsize} shows the Armijo step sizes for gradient ascent.
The step sizes remain small throughout, reflecting the relatively flat landscape in
the 19-dimensional space.
Figure~\ref{fig:cond} shows the Hessian condition number for Newton's method, which
remains large throughout, explaining the suboptimal convergence.

\begin{figure}[H]
  \centering
  \includegraphics[width=0.75\linewidth]{hierarchical_linear_gradient_stepsize.png}
  \caption{Armijo step sizes for gradient ascent per iteration (log scale).}
  \label{fig:stepsize}
\end{figure}

\begin{figure}[H]
  \centering
  \includegraphics[width=0.65\linewidth]{hierarchical_linear_newton_condition.png}
  \caption{Hessian condition number at each Newton iteration.
    The large condition numbers throughout indicate ill-conditioning, consistent with
    Newton's failure to find the CAVI optimum.}
  \label{fig:cond}
\end{figure}

\subsection{Gibbs Sampler Diagnostics}

Figures~\ref{fig:trace_b0}--\ref{fig:trace_tauu} show trace plots for the four
scalar parameters monitored ($\beta_0$, $\beta_1$, $\tau_e$, $\tau_u$).
Figure~\ref{fig:acf} shows the ACF for each parameter; all drop inside the 95\%
confidence band within 5--10 lags, indicating good mixing.
Figure~\ref{fig:rhat} shows the Gelman--Rubin $\hat{R}$ diagnostic; all values
are below 1.01, confirming convergence of the 3 chains.

\begin{figure}[H]
  \centering
  \includegraphics[width=\linewidth]{hierarchical_linear_gibbs_trace_beta0.png}
  \caption{Gibbs trace: $\beta_0$. Red dashed: true value.}
  \label{fig:trace_b0}
\end{figure}

\begin{figure}[H]
  \centering
  \includegraphics[width=\linewidth]{hierarchical_linear_gibbs_trace_beta1.png}
  \caption{Gibbs trace: $\beta_1$.}
  \label{fig:trace_b1}
\end{figure}

\begin{figure}[H]
  \centering
  \includegraphics[width=\linewidth]{hierarchical_linear_gibbs_trace_tau_e.png}
  \caption{Gibbs trace: $\tau_e$ (noise precision).}
  \label{fig:trace_taue}
\end{figure}

\begin{figure}[H]
  \centering
  \includegraphics[width=\linewidth]{hierarchical_linear_gibbs_trace_tau_u.png}
  \caption{Gibbs trace: $\tau_u$ (random-effect precision).}
  \label{fig:trace_tauu}
\end{figure}

\begin{figure}[H]
  \centering
  \includegraphics[width=0.95\linewidth]{hierarchical_linear_gibbs_acf_beta0.png}
  \caption{ACF for $\beta_0$ (chain 1). Red dashed lines: 95\% band.}
  \label{fig:acf}
\end{figure}

\begin{figure}[H]
  \centering
  \includegraphics[width=0.95\linewidth]{hierarchical_linear_gibbs_acf_beta1.png}
  \caption{ACF for $\beta_1$.}
\end{figure}

\begin{figure}[H]
  \centering
  \includegraphics[width=0.95\linewidth]{hierarchical_linear_gibbs_acf_tau_e.png}
  \caption{ACF for $\tau_e$.}
\end{figure}

\begin{figure}[H]
  \centering
  \includegraphics[width=0.95\linewidth]{hierarchical_linear_gibbs_acf_tau_u.png}
  \caption{ACF for $\tau_u$.}
\end{figure}

\begin{figure}[H]
  \centering
  \includegraphics[width=0.55\linewidth]{hierarchical_linear_gibbs_rhat.png}
  \caption{Gelman--Rubin $\hat{R}$ diagnostic for all monitored parameters.
    All values below 1.01 confirm convergence.}
  \label{fig:rhat}
\end{figure}

\subsection{Random Effects Recovery}

Figure~\ref{fig:re} compares the CAVI posterior means for $u_j$ against the
Gibbs reference and the true values.
Both methods recover the group-specific shifts, with CAVI showing slightly tighter
credible intervals consistent with its general tendency towards variance collapse.

\begin{figure}[H]
  \centering
  \includegraphics[width=\linewidth]{hierarchical_linear_random_effects.png}
  \caption{Random effect estimates $u_j$ for CAVI (left) and Gibbs (right).
    Stars indicate true values; error bars show $\pm 2$ posterior SD.}
  \label{fig:re}
\end{figure}

\subsection{Posterior Comparison: VI versus Gibbs}

Figures~\ref{fig:vbgibbs_b0}--\ref{fig:vbgibbs_tauu} overlay the VI approximations
against the Gibbs histograms for the four scalar parameters.

\begin{figure}[H]
  \centering
  \includegraphics[width=0.80\linewidth]{hierarchical_linear_vb_gibbs_beta0.png}
  \caption{Posterior comparison for $\beta_0$: Gibbs histogram versus VI Gaussian
    approximations. The VI curves are dramatically narrower than Gibbs (collapse).}
  \label{fig:vbgibbs_b0}
\end{figure}

\begin{figure}[H]
  \centering
  \includegraphics[width=0.80\linewidth]{hierarchical_linear_vb_gibbs_beta1.png}
  \caption{Posterior comparison for $\beta_1$.
    VI and Gibbs agree well; minimal collapse.}
  \label{fig:vbgibbs_b1}
\end{figure}

\begin{figure}[H]
  \centering
  \includegraphics[width=0.80\linewidth]{hierarchical_linear_vb_gibbs_tau_e.png}
  \caption{Posterior comparison for $\tau_e$.}
  \label{fig:vbgibbs_taue}
\end{figure}

\begin{figure}[H]
  \centering
  \includegraphics[width=0.80\linewidth]{hierarchical_linear_vb_gibbs_tau_u.png}
  \caption{Posterior comparison for $\tau_u$.}
  \label{fig:vbgibbs_tauu}
\end{figure}

\subsection{Posterior Standard-Deviation Ratios}

Figure~\ref{fig:sdratios} shows the ratio $\sigma_{\mathrm{VI}}/\sigma_{\mathrm{Gibbs}}$
for each method and parameter.
CAVI and BFGS exhibit severe collapse for $\beta_0$
(SD ratio $\approx 0.13$) while maintaining near-unity ratios for $\beta_1$.
This asymmetric collapse arises because $\beta_0$ (the intercept) is confounded with the
group means $u_j$ under the mean-field factorisation: the posterior covariance between
$\beta_0$ and $\mathbf{u}$ is discarded, causing $q(\beta_0)$ to appear far more
concentrated than the true marginal.
Newton's method produces anomalous SD ratios ($\beta_1 \approx 2.8$) due to its
convergence to a suboptimal point.

\begin{figure}[H]
  \centering
  \includegraphics[width=0.85\linewidth]{hierarchical_linear_sd_ratios.png}
  \caption{Posterior SD ratios (VI / Gibbs). Dashed line: ideal ratio 1.0.
    Severe intercept collapse ($\approx 0.13$) for CAVI, gradient ascent, and BFGS.}
  \label{fig:sdratios}
\end{figure}

\subsection{Posterior Mean Errors}

Figure~\ref{fig:meanerr} shows the absolute error $|\mu_{\mathrm{VI}} - \mu_{\mathrm{Gibbs}}|$.
Despite severe variance collapse for $\beta_0$, posterior means are well recovered by
all converged methods.

\begin{figure}[H]
  \centering
  \includegraphics[width=0.85\linewidth]{hierarchical_linear_mean_errors.png}
  \caption{Absolute posterior mean errors relative to the Gibbs reference.
    Posterior means are accurate even where variance is collapsed.}
  \label{fig:meanerr}
\end{figure}

\subsection{Performance Summary}

Table~\ref{tab:perf} summarises the convergence, runtime, and final ELBO for all methods.
Table~\ref{tab:post} summarises the posterior means and SD ratios.

\begin{table}[H]
  \centering
  \caption{Optimisation performance: hierarchical linear regression.}
  \label{tab:perf}
  \begin{tblr}{colspec={lrrrr}, hline{1,2,Z}={solid}}
  Method & Iterations & Runtime (ms) & Speedup vs Gibbs & Final ELBO \\
  CAVI & 102 & 27.0 & 121x & -132.077 \\
  Gradient Ascent & 2000 & 8910.5 & 0x & -132.086 \\
  Newton's Method & 50 & 3316.8 & 1.0x & -165.282 \\
  BFGS & 34 & 123.4 & 26x & -132.077 \\
  Gibbs (3 chains) & 3x5000 & 3254.6 & 1x (reference) & N/A \\
\end{tblr}

\end{table}

\begin{table}[H]
  \centering
  \caption{Posterior summary: hierarchical linear regression.}
  \label{tab:post}
  \begin{tblr}{colspec={lrrrrrr}, hline{1,2,Z}={solid}}
  Method & $\mu_{\beta_0}$ & $\mu_{\beta_1}$ & $\mathrm{E}[\tau_e]$ & $\mathrm{E}[\tau_u]$ & SD ratio $\beta_0$ & SD ratio $\beta_1$ \\
  CAVI & 0.4604 & 1.9877 & 3.6959 & 3.2818 & 0.130 & 0.987 \\
  Gradient Ascent & 0.4602 & 1.9878 & 3.6945 & 3.2895 & 0.130 & 0.988 \\
  Newton's Method & 0.4280 & 1.9959 & 2.8896 & 1.9989 & 0.325 & 2.779 \\
  BFGS & 0.4603 & 1.9877 & 3.6959 & 3.2818 & 0.130 & 0.987 \\
  Gibbs (ref) & 0.4257 & 1.9872 & 3.7214 & 2.8239 & 1.000 & 1.000 \\
\end{tblr}

\end{table}

\subsection{Computational Cost}

Figure~\ref{fig:timing} compares the runtimes of all methods.
CAVI is the fastest converged method at 27\,ms, 134$\times$ faster than the
conjugate Gibbs sampler (3.6\,s for 3 chains $\times$ 5000 iterations).
BFGS is 31$\times$ faster than Gibbs.
Gradient ascent reaches the same final ELBO as CAVI but requires the full 2000-iteration
budget, making it $\sim$\,3$\times$ slower than Gibbs at this problem size.
Newton's method, despite only 50 iterations, requires a $19\times19$ FD Hessian at each
step (requiring $19^2 = 361$ ELBO evaluations per iteration), giving a total cost
comparable to Gibbs.

\begin{figure}[H]
  \centering
  \includegraphics[width=0.85\linewidth]{hierarchical_linear_timing_comparison.png}
  \caption{Runtime comparison (log scale). Annotated speedup factors relative to Gibbs.}
  \label{fig:timing}
\end{figure}

%─────────────────────────────────────────────────────────────────────────────
\section{Discussion}
%─────────────────────────────────────────────────────────────────────────────

\paragraph{Posterior variance collapse.}
The intercept $\beta_0$ exhibits severe collapse (SD ratio $\approx 0.13$) under all
well-converged VI methods.
The mechanism is the mean-field assumption: the true posterior has strong negative
correlation between $\beta_0$ and the group means $\bar{u} = J^{-1}\sum_j u_j$,
because the data constrain $\beta_0 + \bar{u}$ far more tightly than either component
alone.
Mean-field VI factorises these two blocks, forcing $q(\beta_0)$ to be a narrow Gaussian
centred at the marginal mode rather than the true marginal.
The slope $\beta_1$ is not confounded with the random effects in the same way, so its
VI standard deviation matches Gibbs closely.

\paragraph{Newton's method failure.}
Newton's method converges to a point with ELBO $= -165.28$, substantially worse than
the global CAVI optimum at $-132.08$.
This is a consequence of the finite-difference $19\times19$ Hessian having poor
conditioning throughout the trajectory: the Hessian of the ELBO mixes parameters
that operate on very different scales ($\bm{\mu}_\beta$ in $[-2,2]$, $\log b_u$ near
$[-3, 3]$), making the Newton step unreliable even with regularisation.
BFGS avoids this issue by building an \emph{adaptive} positive-definite inverse-Hessian
approximation, which effectively rescales each parameter direction from curvature
information accumulated over iterations.

\paragraph{CAVI versus BFGS.}
CAVI and BFGS converge to the same ELBO optimum with comparable accuracy.
CAVI is $\sim 4\times$ faster in this implementation because each iteration requires
only closed-form matrix operations, whereas BFGS requires $2D = 38$ ELBO evaluations
per iteration for FD gradients.
For models where CAVI updates are not available (e.g.\ non-conjugate likelihoods),
BFGS is the natural first-choice gradient-based method.

\paragraph{Gradient ascent.}
First-order gradient ascent converges to the correct optimum but requires all 2000
iterations, making it the slowest converged method.
The Armijo step sizes (Figure~\ref{fig:stepsize}) show consistently small values,
indicating that the curvature of the ELBO surface is poorly aligned with the gradient
direction in a 19-dimensional space --- the same limitation identified in the simpler
linear and quadratic case studies, now more pronounced.

%─────────────────────────────────────────────────────────────────────────────
\section{Conclusion}
%─────────────────────────────────────────────────────────────────────────────

Hierarchical Bayesian linear regression with $J=5$ groups and $D=19$ unconstrained
parameters presents a qualitatively more challenging optimisation problem than the
two-parameter logistic case.
CAVI and BFGS both find the global mean-field ELBO optimum and recover accurate
posterior means, but both exhibit severe variance collapse for the intercept $\beta_0$
due to the mean-field factorisation ignoring the posterior correlation between the
intercept and the random effects.
Newton's method fails to converge to this optimum, highlighting the limitations of
finite-difference second-order methods when the Hessian is ill-conditioned.
Gradient ascent converges but is an order of magnitude slower.
The conjugate Gibbs sampler provides accurate marginal posteriors but at a cost
$>100\times$ that of CAVI for this problem size.

\printbibliography

\end{document}
